\section*{Abstract}
\thispagestyle{empty}

The project titled “Development of an Online Library Management System” was conducted to fulfill digital transformation demands in library resource management and utilization. In the context of traditional libraries gradually transitioning to digital environments, constructing a modern, efficient management system capable of providing better support for users is an imperative requirement. The project focuses on researching, designing, and deploying a library system as a Web Application targeted at small-to-medium scale libraries, such as school libraries, faculty libraries, or learning resource centers. \\

The system is constructed to manage diverse types of library resources, including printed books, journals, and digital assets such as E-books and PDF documents. Furthermore, the system goes beyond supporting basic administrative functions by aiming to enhance user experience through the integration of modern technologies, particularly Artificial Intelligence (AI). \\

Within the system, the primary user roles consist of readers, librarians, and administrators. Each role possesses distinct functional requirements and usage patterns, which are appropriately supported by the system.\\

For the reader role, the system provides flexible document search and retrieval functions through keywords or intelligent search tools. Additionally, the system incorporates a document recommendation mechanism based on reading history and user behavior, helping readers quickly access documents tailored to their learning and research needs. Readers can manage their borrowing history, track loan-return statuses, and interact with the system via a chatbot for prompt assistance with queries.\\

For the librarian role, the system supports managing resource catalogs, updating document metadata, and processing borrowing and returning operations. Furthermore, librarians can monitor resource availability, control data, and ensure efficient library operations. Implementing the system minimizes manual procedures, thereby increasing productivity and reducing management errors.\\

For the administrator role, the system provides management privileges over all system operations, including user account management, system configuration, and data monitoring. In addition, administrators can review statistics related to borrowing and returning activities to support managerial decision-making and rational library resource development.\\

A key highlight of the system is the integration of Artificial Intelligence capabilities, such as document recommendations, natural language search, and an interactive assistant chatbot. These features significantly enhance user-system interaction and elevate the overall user experience compared to traditional library systems.\\

Architecturally, the system is designed following a web application model with key components: user interface, server-side business processing logic, and data storage databases. This organization ensures flexibility, scalability, and maintainability for future upgrades.\\

Throughout the development of the Online Library Management System, this project aims to optimize resource management workflows, enhance operational efficiency, and deliver an improved user experience in searching for and utilizing academic materials. Simultaneously, the project demonstrates the practical applicability of modern technologies, especially AI, in developing information systems for educational and research environments.\\

\clearpage
\pagenumbering{arabic}
